\chapter{Versuchsaufbau und Evaluationsmethodik}
\label{chap:versuchsaufbau-evaluationsmethodik}

\section{Versuchsplan und Messkonfiguration}
\label{sec:versuchsplan}
Um reproduzierbare Ergebnisse zu erzielen, folgt jeder Störungsversuch einem standardisierten Ablauf:
\begin{enumerate}
    \item Herstellung eines definierten Ausgangszustands (Baseline) mit bekannter Commit-ID und Image-Version.
    \item Start der kontinuierlichen Hintergrundlast (k6-Lasttest).
    \item Gezielte Induktion der Störung zum definierten Zeitpunkt $t_{\text{fault}}$.
    \item Kontinuierliche Aufzeichnung von Fehlerraten und Antwortzeiten bis zur vollständigen Stabilisierung $t_{\text{recovered}}$.
    \item Protokollierung und Export der Rohdaten.
\end{enumerate}

\section{Basiskonfiguration und verbesserte Konfiguration}
\label{sec:basis-vs-verbessert}
Für die Experimente werden jeweils zwei Szenarien verglichen:
\begin{itemize}
    \item \textbf{Basiskonfiguration (Single Replica):} Minimalauslegung mit exakt einer Pod-Instanz ohne PodDisruptionBudget.
    \item \textbf{Verbesserte Konfiguration (Multi Replica):} Redundante Auslegung mit mindestens zwei Pod-Instanzen, feingranularen Health Probes und Rolling-Update-Strategie.
\end{itemize}

\section{Lastprofil und Messwerkzeuge}
\label{sec:lastprofil-werkzeuge}
Zur Lastgenerierung wird das Open-Source-Tool k6 eingesetzt. Ein konstanter Anfragestrom von 10 bis 20 virtuellen Benutzern (VUs) fragt den Statusendpunkt ab. Gemessen werden HTTP-Statuscodes, Fehlerraten sowie die Latenzperzentile (Median p50 und 95.\,Perzentil p95).

\section{Definition von Ausfall und Wiederherstellung}
\label{sec:definition-ausfall-recovery}
Ein Systemzustand gilt als \textbf{ausgefallen}, sobald die HTTP-Erfolgsrate (Statuscode 200) unter 99\,\% fällt oder Timeouts auftreten. Die \textbf{Wiederherstellungszeit} (Mean Time to Recovery, MTTR) bezeichnet die Zeitspanne vom Auftreten des ersten Fehlers bis zum Erreichen eines stabilen Betriebs ohne Fehlantworten über ein Intervall von mindestens 30 Sekunden.
